% GENERATED FILE. Edit canonical lessons, then run npm run generate:books.
% canonical-source-hash: fnv1a64:3938f6eba10beeca
% canonical-lessons: FA-C291-tarif-kardan, FA-W291-eyn, FA-C291-avval-bad, FA-C291-vaqti, FA-C291-khandani, FA-R291-a-story-in-order

\chapter{Then, And In The End --- A Story In Order}
\label{ch:fa-a-story-in-order}
\hlchaptermodality{291}

\begin{chapteropening}
\textbf{By the end of this chapter:} \emph{I can tell a short story in order with avval, ba'd and dar akhar, join events with vaqti, and read a connected story of about a hundred words.}

Complete FA-R291-a-story-in-order: retell a short trip in order, with a scene, two events and an ending.
\end{chapteropening}

\section[\texorpdfstring{\fa{تعریف} \fa{کردن} (ta'rif kardan)}{تعریف کردن (ta'rif kardan)}]{\fa{تعریف} \fa{کردن} (ta'rif kardan) — to tell what happened}

\label{lesson:FA-C291-tarif-kardan}



\begin{quote}
Say \textquotedblleft{}a story\textquotedblright{}, then \textquotedblleft{}I was reading\textquotedblright{}. (\emph{Dâstân}, \emph{mikhândam}.)

Chapter two hundred and ninety gave you the two pasts a story is made of. This is the verb for telling one.
\end{quote}

\subsection*{You'll want to know: \fa{تعریف} \fa{کردن}}
\textbf{\fa{تعریف} \fa{کردن}} — \emph{ta'rif kardan} — \textquotedblleft{}to tell\textquotedblright{}, of a story or of what happened.

\begin{quote}
\textbf{\fa{او} \fa{داستان} \fa{سفر} \fa{را} \fa{تعریف} \fa{کرد}.} — \emph{u dâstân-e safar râ ta'rif kard} — She told the story of the trip.
\end{quote}

It runs on \emph{kardan}, so the tenses come for free: \emph{ta'rif mikonam}, I tell; \emph{ta'rif kardam}, I told; \emph{ta'rif mikardam}, I was telling.

Put \emph{az}, \textquotedblleft{}from\textquotedblright{}, in front of a thing and the verb turns to praise: \textbf{\fa{از} \fa{هتل} \fa{تعریف} \fa{کرد}}, \emph{az hotel ta'rif kard}, \textquotedblleft{}he spoke well of the hotel\textquotedblright{}.

Say it \emph{ta'-rif}, with a small catch after \emph{ta}, as if the word drew breath. That catch is a letter of its own, and the next lesson puts it in your hand.

\begin{cousinweb}
\emph{ta'rif} is borrowed from Arabic \emph{taʿrīf} (EtymDB 2.1). It is usually explained as \textquotedblleft{}making known\textquotedblright{}, a noun on the root \emph{ʿ-r-f} of Arabic \emph{ʿarafa}, \textquotedblleft{}to know\textquotedblright{}. \emph{Kardan} is the Persian verb \textquotedblleft{}to do\textquotedblright{} from chapter seven. So \emph{ta'rif kardan} is \textquotedblleft{}to do a making-known\textquotedblright{}: you tell a story, and now the listener knows it. To praise a hotel is to make its good side known.
\end{cousinweb}

\subsection*{Your turn}
Say these aloud:

\begin{itemize}
\raggedright
  \item \emph{ta'rif kardan}, with the catch — \emph{ta'-rif}
  \item \textquotedblleft{}She told the story of the trip\textquotedblright{} — \emph{u dâstân-e safar râ ta'rif kard}
  \item \textquotedblleft{}I tell\textquotedblright{}, \textquotedblleft{}I told\textquotedblright{}, \textquotedblleft{}I was telling\textquotedblright{} — \emph{ta'rif mikonam}, \emph{ta'rif kardam}, \emph{ta'rif mikardam}
  \item \textquotedblleft{}He spoke well of the hotel\textquotedblright{} — \emph{az hotel ta'rif kard}
\end{itemize}

\begin{tcolorbox}[breakable,colback=teal!4,colframe=teal!35!black,title={Before you move on}]
\textquotedblleft{}To tell a story\textquotedblright{}? (\emph{\textbf{Dâstân ta'rif kardan.}}) What changes when \emph{az} comes in front? (It becomes \textbf{to speak well of}.) Which Arabic root is inside \emph{ta'rif}? (\emph{\textbf{ʿ-r-f}}, \textquotedblleft{}to know\textquotedblright{}.)
\end{tcolorbox}
\section[\texorpdfstring{\fa{ع} ('eyn)}{ع ('eyn)}]{\fa{ع} ('eyn) — the catch inside \fa{تعریف}}

\label{lesson:FA-W291-eyn}



\begin{quote}
Before the new one: write \fa{ت} and \fa{ر} once each, and name them.

Now say \textbf{\fa{تعریف} \fa{کردن}} (\emph{ta'rif kardan}), \textquotedblleft{}to tell\textquotedblright{}. You met it a moment ago. This lesson takes the catch out of it and puts it in your hand.
\end{quote}

\subsection*{Script you'll notice: \fa{ع}}
\textbf{\fa{ع}} — \emph{'eyn}, the small catch written \emph{'} in this book's romanization.

It is the second letter of \textbf{\fa{تعریف}}, straight after \textbf{\fa{ت}}. In Persian it is a light stop in the throat, \emph{ta'-rif}, or no more than a slightly longer vowel before it. Arabic says it deeper in the throat; Persian does not.

It changes shape more than most letters. At the start of a word it is an open hook, \textbf{\fa{عـ}}. In the middle it closes into a small flat loop, \textbf{\fa{ـعـ}}, which is the shape inside \textbf{\fa{تعریف}}. Standing alone, \textbf{\fa{ع}}, the hook sits on a tail that curls down below the line.

Almost every Persian word with \textbf{\fa{ع}} came from Arabic, and Persian has thousands of them.

\subsection*{Writing: \fa{ع}}
Put your pen on \fa{ع} and follow its line. Copy the shape you can see — slowly, and larger than it is printed.

\begin{quote}
This book does not yet tell you \textbf{where to start this letter or which way to travel}. It is not written down here until it can be written down with a source. Copying what is in front of you needs no such source, and it is how the shape gets into your hand in the meantime.
\end{quote}

\subsection*{Your turn}
\begin{itemize}
\raggedright
  \item \emph{Find:} the closed \fa{ع} inside \textbf{\fa{تعریف}}, and point at it
  \item \emph{Trace it:} \fa{ع} standing alone three times, saying the catch as you finish each one
  \item \emph{Write it:} \fa{ت}, then \fa{تع}, joined, from memory
  \item \emph{Write it:} \textbf{\fa{تعریف}}
  \item \emph{Read it:} \textbf{\fa{تعریف} \fa{کردن}} aloud once more, with the catch
\end{itemize}

\begin{tcolorbox}[breakable,colback=teal!4,colframe=teal!35!black,title={Before you move on}]
Which letter is this — \fa{ع}? (\emph{\textbf{'eyn}}, a light catch.) Which word did it come from? (\textbf{\fa{تعریف} \fa{کردن}}, \emph{ta'rif kardan}, \textquotedblleft{}to tell\textquotedblright{}.) Which shape does it take in the middle of a word? (\textbf{A small closed loop.})
\end{tcolorbox}
\section[\texorpdfstring{\fa{اول،} \fa{بعد،} \fa{در} \fa{آخر} (avval, ba'd, dar …)}{اول، بعد، در آخر (avval, ba'd, dar …)}]{\fa{اول،} \fa{بعد،} \fa{در} \fa{آخر} (avval, ba'd, dar âkhar) — first, then, in the end}

\label{lesson:FA-C291-avval-bad}



\begin{quote}
Say \textquotedblleft{}first\textquotedblright{}, then \textquotedblleft{}then, afterwards\textquotedblright{}, then \textquotedblleft{}at last\textquotedblright{}. (\emph{Avval}, \emph{sepas}, \emph{bâlâkhare}.)

You hold the words already. This lesson strings them through a run of events, and adds the one people use most in speech.
\end{quote}

\begin{grammarlens}[title={Grammar Lens: three events, in order}]
\begin{quote}
\textbf{\fa{اول} \fa{صبحانه} \fa{خوردم}. \fa{بعد} \fa{به} \fa{کتابخانه} \fa{رفتم}. \fa{در} \fa{آخر} \fa{به} \fa{خانه} \fa{برگشتم}.}
\end{quote}

\begin{quote}
\emph{avval sobhâne khordam. ba'd be ketâbkhâne raftam. dar âkhar be khâne bargashtam.}
\end{quote}

\begin{quote}
First I had breakfast. Then I went to the library. In the end I went back home.
\end{quote}

\textbf{\fa{بعد}}, \emph{ba'd}, \textquotedblleft{}then\textquotedblright{}, opens the next event. It does the job \emph{sepas} does, and it is the one you hear at a kitchen table; \emph{sepas} sounds written.

With \emph{az} after it, \emph{ba'd az} is \textquotedblleft{}after\textquotedblright{}: \textbf{\fa{بعد} \fa{از} \fa{صبحانه}}, \emph{ba'd az sobhâne}, \textquotedblleft{}after breakfast\textquotedblright{}. \emph{Pas az}, from chapter one hundred and fifty-two, means the same and sounds more formal.

The events stay in the plain past, one after another, each with its edge.
\end{grammarlens}

\begin{grammarlens}[title={Grammar Lens: in the end, and at last}]
\textbf{\fa{در} \fa{آخر}}, \emph{dar âkhar}, is \textquotedblleft{}in the end\textquotedblright{}: \emph{dar}, \textquotedblleft{}in\textquotedblright{}, and \emph{âkhar}, \textquotedblleft{}the end\textquotedblright{}. It closes the run.

\emph{Bâlâkhare}, \textquotedblleft{}at last\textquotedblright{}, closes it too, with relief in it: the thing you were waiting for has come. \emph{Bâlâkhare residam}: I got there at last.
\end{grammarlens}

\begin{cousinweb}
\emph{ba'd} is borrowed from Arabic \emph{baʿd}, \textquotedblleft{}after\textquotedblright{} (EtymDB 2.1), and it carries the catch you wrote one lesson ago: \textbf{\fa{ب}}, then \textbf{\fa{ع}}, then \textbf{\fa{د}}. \emph{Avval} and \emph{âkhar} are Arabic as well (EtymDB 2.1). In these three sentences the words that order the events are Arabic, and the verbs are Persian.
\end{cousinweb}

\subsection*{Your turn}
Say these aloud:

\begin{itemize}
\raggedright
  \item \emph{avval sobhâne khordam. ba'd be ketâbkhâne raftam.}
  \item add the end — \emph{dar âkhar be khâne bargashtam}
  \item \textquotedblleft{}after breakfast\textquotedblright{}, two ways — \emph{ba'd az sobhâne}, \emph{pas az sobhâne}
  \item three things you did yesterday, with \emph{avval}, \emph{ba'd} and \emph{dar âkhar}
\end{itemize}

\begin{tcolorbox}[breakable,colback=teal!4,colframe=teal!35!black,title={Before you move on}]
Which word for \textquotedblleft{}then\textquotedblright{} sounds like speech, and which sounds written? (\emph{\textbf{Ba'd}} for speech, \emph{\textbf{sepas}} for writing.) \textquotedblleft{}After breakfast\textquotedblright{}? (\emph{\textbf{Ba'd az sobhâne.}}) What does \emph{bâlâkhare} add that \emph{dar âkhar} does not? (\textbf{Relief}: at last.)
\end{tcolorbox}
\section[\texorpdfstring{\fa{وقتی} (vaqti)}{وقتی (vaqti)}]{\fa{وقتی} (vaqti) — when, and a story told in order}

\label{lesson:FA-C291-vaqti}



\begin{quote}
Say \textquotedblleft{}first\textquotedblright{}, \textquotedblleft{}then\textquotedblright{}, \textquotedblleft{}in the end\textquotedblright{}. (\emph{Avval}, \emph{ba'd}, \emph{dar âkhar}.) Then ask \textquotedblleft{}when?\textquotedblright{} (\emph{Che vaqt?})
\end{quote}

\begin{grammarlens}[title={Grammar Lens: vaqti, the joint for a moment}]
\textbf{\fa{وقتی}}, \emph{vaqti}, is \textquotedblleft{}when\textquotedblright{} in the middle of a story, not in a question. It ties what follows to a moment:

\begin{quote}
\textbf{\fa{وقتی} \fa{به} \fa{هتل} \fa{رسیدم،} \fa{باد} \fa{میوزید}.} — \emph{vaqti be hotel residam, bâd mivazid} — When I got to the hotel, the wind was blowing.
\end{quote}

\emph{Vaqti} takes the event, \emph{residam}, and the rest paints what was going on at that moment, \emph{mivazid}.

It is \emph{vaqt}, \textquotedblleft{}time\textquotedblright{}, with \emph{-i}: \textquotedblleft{}at the time that\textquotedblright{}. \emph{Vaqt} is the word inside \emph{che vaqt} and \emph{sar-e vaqt}. The letter \textbf{\fa{ق}} is written \emph{q} here, so chapter seventy's \emph{ghatâr}, \textquotedblleft{}a train\textquotedblright{}, appears as \emph{qatâr}.
\end{grammarlens}

\begin{grammarlens}[title={Grammar Lens: bâ, \textquotedblleft{}with\textquotedblright{}}]
\textbf{\fa{با}}, \emph{bâ}, is \textquotedblleft{}with\textquotedblright{}: \emph{bâ yek dust}, with a friend. For travel it is \textquotedblleft{}by\textquotedblright{}: \emph{bâ qatâr}, by train. And \emph{bâ ham} is \textquotedblleft{}together\textquotedblright{}.
\end{grammarlens}

\begin{grammarlens}[title={Grammar Lens: the order a told story keeps}]
A story you tell, that you \emph{ta'rif mikonid}, keeps three parts in one order. First the time and the scene, in \emph{bud} and the \emph{mi-} past: \emph{sâl-e gozashte... havâ sard bud.} Then the events, in the plain past, strung on \emph{avval}, \emph{ba'd} and \emph{vaqti}. Then the close, with \emph{dar âkhar}.

\begin{quote}
\textbf{\fa{سال} \fa{گذشته} \fa{با} \fa{قطار} \fa{به} \fa{تهران} \fa{رفتم}. \fa{وقتی} \fa{رسیدم،} \fa{هوا} \fa{سرد} \fa{بود}. \fa{اول} \fa{به} \fa{هتل} \fa{رفتم}. \fa{بعد} \fa{با} \fa{یک} \fa{دوست} \fa{شام} \fa{خوردم}. \fa{در} \fa{آخر} \fa{به} \fa{هتل} \fa{برگشتم}.}
\end{quote}

\begin{quote}
\emph{sâl-e gozashte bâ qatâr be Tehrân raftam. vaqti residam, havâ sard bud. avval be hotel raftam. ba'd bâ yek dust shâm khordam. dar âkhar be hotel bargashtam.}
\end{quote}

\begin{quote}
Last year I went to Tehran by train. When I arrived, it was cold. First I went to the hotel. Then I had dinner with a friend. In the end I went back to the hotel.
\end{quote}
\end{grammarlens}

\subsection*{Your turn}
Say these aloud:

\begin{itemize}
\raggedright
  \item \emph{vaqti be hotel residam, bâd mivazid}
  \item \textquotedblleft{}by train\textquotedblright{}, then \textquotedblleft{}with a friend\textquotedblright{} — \emph{bâ qatâr}, \emph{bâ yek dust}
  \item the Tehran story, once, without stopping
  \item your own: one line of scene, two events with \emph{avval} and \emph{ba'd}, and a close with \emph{dar âkhar}
\end{itemize}

\begin{tcolorbox}[breakable,colback=teal!4,colframe=teal!35!black,title={Before you move on}]
What does \emph{vaqti} do that \emph{che vaqt} does not? (It joins two parts of a story and asks nothing.) \textquotedblleft{}Together\textquotedblright{}? (\emph{\textbf{Bâ ham.}}) What comes first in a told story? (\textbf{The time and the scene.})
\end{tcolorbox}
\section[(reading)]{(reading) — a winter trip, told in both pasts}

\label{lesson:FA-C291-khandani}



\begin{quote}
Chapter twenty-one gave you six words to read cold. This passage is ninety-three words long, and it happened last winter.

Every word in it was taught before this lesson.
\end{quote}

\begin{tcolorbox}[breakable,skin=enhanced,colback=orange!6,colframe=orange!45!black,boxrule=0.5pt,arc=1mm,left=6pt,right=6pt,top=4pt,bottom=4pt,fonttitle=\bfseries,title={Read this}]
\begin{quote}
\fa{سال} \fa{گذشته،} \fa{در} \fa{زمستان،} \fa{با} \fa{یک} \fa{دوست} \fa{به} \fa{تهران} \fa{رفتم}. \fa{هوا} \fa{خیلی} \fa{سرد} \fa{بود} \fa{و} \fa{باد} \fa{میوزید}. \fa{وقتی} \fa{رسیدیم،} \fa{اول} \fa{به} \fa{هتل} \fa{رفتیم}. \fa{بعد} \fa{در} \fa{یک} \fa{رستوران} \fa{قدیمی} \fa{شام} \fa{خوردیم}. \fa{رستوران} \fa{گرم} \fa{و} \fa{آرام} \fa{بود}. \fa{در} \fa{حالی} \fa{که} \fa{شام} \fa{میخوردیم،} \fa{او} \fa{یک} \fa{داستان} \fa{جالب} \fa{تعریف} \fa{میکرد} \fa{و} \fa{من} \fa{گوش} \fa{میدادم}.
\end{quote}

\begin{quote}
\fa{صبح} \fa{به} \fa{موزه} \fa{رفتیم} \fa{و} \fa{در} \fa{پارک} \fa{قدم} \fa{زدیم}. \fa{خورشید} \fa{میتابید،} \fa{ولی} \fa{هنوز} \fa{سرد} \fa{بود}. \fa{بعد} \fa{از} \fa{ناهار} \fa{به} \fa{بازار} \fa{رفتیم} \fa{و} \fa{یک} \fa{کیف} \fa{خریدیم}. \fa{در} \fa{آخر} \fa{با} \fa{قطار} \fa{به} \fa{خانه} \fa{برگشتیم}. \fa{سفر} \fa{کوتاه} \fa{بود،} \fa{ولی} \fa{خیلی} \fa{خوب} \fa{بود}. \fa{دوباره} \fa{به} \fa{تهران} \fa{میروم}.
\end{quote}

Read it once without stopping. Do not translate as you go; let the sentences arrive.

Then read it a second time and sort every verb into one of two piles: the scene, or what happened.
\end{tcolorbox}

\begin{culture}
The events: \emph{raftam}, \emph{residim}, \emph{raftim}, \emph{khordim}, \emph{qadam zadim}, \emph{kharidim}, \emph{bargashtim}. Each one happened and was over, and \emph{avval}, \emph{ba'd} and \emph{dar âkhar} keep them in order.

The scene: the cold, the wind, the warm quiet restaurant, the sun. \emph{Bud} paints it, and so does the \emph{mi-} past: \emph{mivazid}, \emph{mitâbid}.

At dinner both verbs are scene: \emph{ta'rif mikard} and \emph{gush midâdam} ran side by side, with \emph{dar hâli ke} in front. The last line is in the present, \emph{miravam}, because it is a plan.
\end{culture}

\subsection*{Your turn}
\begin{itemize}
\raggedright
  \item \emph{Say it:} the passage aloud, once, without stopping
\end{itemize}

\begin{tcolorbox}[breakable,colback=teal!4,colframe=teal!35!black,title={Before you move on}]
How many words in that passage were new? (\textbf{None.})

Why is \emph{mivazid} in the \emph{mi-} past? (The wind \textbf{was blowing}: it is scene.)

Which word opened the line where they arrived? (\emph{\textbf{Vaqti.}})
\end{tcolorbox}
\section[(review)]{Review — telling the trip back}

\label{lesson:FA-R291-a-story-in-order}



\begin{quote}
Say \textquotedblleft{}she was telling a story\textquotedblright{}. (\emph{U dâstân ta'rif mikard.}) Then \textquotedblleft{}I was listening\textquotedblright{}. (\emph{Gush midâdam.})
\end{quote}

\begin{grammarlens}[title={What you've built}]
Two chapters, and a whole story fits in your mouth. The time first, with \emph{gozashte}. The scene in \emph{bud} and the \emph{mi-} past. The events in the plain past, strung on \emph{avval}, \emph{ba'd} and \emph{vaqti}. The close with \emph{dar âkhar}. \emph{Bâ} tells you who came along and how you travelled, and \emph{ta'rif kardan} is what you are doing while you say it.

Now tell the winter trip back, not word for word, and in the first person plural, \emph{mâ}. The skeleton:

\begin{quote}
\emph{sâl-e gozashte be Tehrân raftim... havâ sard bud... vaqti residim... avval... ba'd... dar âkhar bâ qatâr bargashtim.}
\end{quote}
\end{grammarlens}

\subsection*{Your turn}
Say these aloud:

\begin{itemize}
\raggedright
  \item the skeleton, filled in, in at least five sentences
  \item one thing they did not do — \emph{be sinemâ naraftim}, they didn't go to the cinema
  \item \textquotedblleft{}When we arrived, it was cold\textquotedblright{} — \emph{vaqti residim, havâ sard bud}
  \item \textquotedblleft{}Last week I told the story\textquotedblright{} — \emph{hafte-ye gozashte dâstân râ ta'rif kardam}
\end{itemize}

\begin{tcolorbox}[breakable,colback=teal!4,colframe=teal!35!black,title={Before you move on}]
What are the three parts of a told story, in order? (\textbf{The time and scene, the events, the close.}) Which joint ties an event to its moment? (\emph{\textbf{Vaqti.}}) Which letter in \textbf{\fa{تعریف}} is the catch? (\textbf{\fa{ع}}, \emph{'eyn}.) And which one in \textbf{\fa{گذشته}} says \emph{z}? (\textbf{\fa{ذ}}, \emph{zâl}.)
\end{tcolorbox}
